\section{Safety Guarantee and Certified Fallback}
\label{section_safety_fallback}

    Although the constrained searching methods in Sections~\ref{section_priority_opt}--\ref{section_config_opt}
    can typically keep important tasks schedulable, Two gaps remain.
    First, the execution-time (ET) prediction the gate is certified against may be inaccurate:
    a gate-passing configuration could still miss deadlines if the environment changes abruptly.
    Second, the search is heuristic and cannot guarantee to find a feasible configuration
    in every environment.

    To close these gaps, this section introduces a \emph{certified safe fallback}---
    a configuration whose schedulability is established offline based on the worst-case scenario analysis to ensure the system can operate safely online if execution time prediction is accurate.

    \subsection{When the Safety Fallback is Deployed}
    \label{section_fallback_triggers}
    At every interval the scheduler holds a certified safe fallback
    $(\boldsymbol{P}^{\star},\boldsymbol{Q}^{\star})$ and an incumbent champion scored against it.
    A surviving candidate is \emph{adopted}---promoted to the deployed configuration---
    only if it strictly improves \eqref{eq_overall_obj};
    the feasibility check is already enforced in-search,
    so adoption never sacrifices the important-task constraint.
    When no strictly-improving feasible candidate is found, the fallback is shipped
    (the deployed configuration is left unchanged).
    The fallback is deployed whenever:
    \begin{itemize}
        \item \textbf{No improving feasible configuration is found} this interval;
        \item \textbf{The inherited incumbent is no longer schedulable} under the updated
        distributions---then the fallback also seeds the search
        (Section~\ref{section_initial_solution});
        \item \textbf{An execution-time change exceeds the drift threshold},
        invalidating the inherited config (Section~\ref{section_initial_solution}).
    \end{itemize}
    Optimization and safety thereby decouple:
    aggressive candidates are explored freely,
    since infeasible or non-improving ones fall back to the certified config.

    \subsection{Constructing the Certified Fallback Offline}
    \label{section_fallback_construction}
    The fallback is constructed offline against a worst-case execution-time (WCET) model,
    so its schedulability is certified independently of the online predictions.
    We assume the task set is schedulable under worst-case execution times;
    otherwise no online algorithm can give a schedulability guarantee.
    Under this assumption the fallback is the \emph{DM-Fast} configuration---
    the most schedulable configuration, and the correct seed for both the fallback
    and the online search.

    \begin{definition}[DM-Fast configuration]
    \label{def_dm_fast}
    The \emph{DM-Fast} configuration assigns each task the smallest QoS budget on its grid
    (the fastest, i.e.\ shortest, execution time)
    and orders tasks by Deadline-Monotonic (DM) priorities---shorter relative deadlines first.
    \end{definition}

    \begin{theorem}[DM-Fast schedulability dominance]
    \label{thm:dm_fast_dominance}
    Under preemptive partitioned scheduling with constrained deadlines $D_i \le T_i$,
    if DM-Fast is unschedulable for some task $\tau_i$,
    then no configuration $(\boldsymbol{P},\boldsymbol{Q})$ is schedulable for $\tau_i$.
    \end{theorem}
    \noindent\textit{Proof.}
    Schedulability of $\tau_i$ depends only on its higher-priority set $\text{hp}(i)$
    and the execution times therein;
    lower-priority and cross-core tasks do not affect $\rtDist{i}$
    (Lemmas~\ref{lem:rta_cross_core}--\ref{lem:rta_lp_pure_move}),
    so the argument is per-core.
    \emph{Priority.} For constrained-deadline tasks, DM is optimal among fixed-priority
    assignments~\cite{Leung1982DeadlineMonotonic}:
    for fixed $\boldsymbol{Q}$, DM minimizes every response time.
    Hence DM-unschedulable under $\boldsymbol{Q}$ implies unschedulable
    under any $\boldsymbol{P}$ with that $\boldsymbol{Q}$.
    \emph{QoS budget.} A larger budget yields a longer execution time
    (Section~\ref{section_rta_distribution}),
    and interference is monotonic in execution time,
    so it can only \emph{increase} response times.
    Hence for fixed $\boldsymbol{P}$, the smallest budget minimizes every response time;
    unschedulable under the smallest budget implies unschedulable under any larger budget.
    Combining both, DM-Fast minimizes every response time along both axes.
    If DM-Fast leaves $\tau_i$ unschedulable,
    no $(\boldsymbol{P},\boldsymbol{Q})$ can schedule it. \hfill$\square$

    By Theorem~\ref{thm:dm_fast_dominance}, the WCET-schedulability assumption is equivalent
    to DM-Fast being schedulable under WCET.
    We therefore take DM-Fast-under-WCET as the certified fallback,
    verified via \eqref{eq_prob_rta_in_opt} to satisfy
    $Pr(r_i > D_i) \leq \Theta_i$ for every $\tau_i \in \boldsymbol{\tau}^{safe}$.
    The same configuration seeds the online search
    (Section~\ref{section_initial_solution});
    since the seed is schedulable, the search only moves between schedulable configurations.
    If even DM-Fast under the predicted distributions is unschedulable for an important task,
    the environment is flagged unschedulable and the previous fallback is retained.

    \subsection{Schedulability Guarantee}
    \label{section_schedulability_guarantee}
    In-search gating plus the certified fallback yield the following guarantee
    on the shipped configuration.

    \begin{theorem}[Important-task schedulability]
    \label{obs:safety_fallback}
    If the execution-time prediction is accurate---predicted distribution equals true---
    then at every shipped interval the deployed configuration
    $(\boldsymbol{P}^{\star},\boldsymbol{Q}^{\star})$ satisfies
    \[
        Pr(r_i > D_i) \leq \Theta_i, \quad \forall \tau_i \in \boldsymbol{\tau}^{safe},
    \]
    regardless of whether the ongoing search has converged.
    \end{theorem}
    \noindent\textit{Proof.}
    The deployed configuration is either (i) a candidate promoted by the adopt-or-retain rule
    (Section~\ref{section_fallback_triggers}),
    or (ii) the certified fallback retained when none is found.
    Case (i): the in-search gates (Algorithms~\ref{alg:modified_audsley}, \ref{alg:qos_walk})
    reject any infeasible candidate before adoption,
    so a promoted candidate satisfies \eqref{eq_important_task_constraint}
    under the predicted distribution;
    with accurate prediction it holds under the true distribution.
    Case (ii): the fallback is WCET-verified to satisfy the bound,
    and true execution time is no worse than WCET.
    In both cases the shipped config satisfies the bound;
    an incomplete optimization may degrade performance but never safety. \hfill$\square$

    This guarantee needs no conservative WCET upper bound for the online search:
    the in-search check is certified against the probabilistic RTA,
    so the online guarantee is tight rather than pessimistic,
    with offline WCET reserved for the fallback alone.
